% !TeX program = pdflatex
% !BIB program = biber
\documentclass[11pt,a4paper]{article}

% --- Idioma y codificación -------------------------------------------------
\usepackage[utf8]{inputenc}
\usepackage[T1]{fontenc}
\usepackage[spanish,es-tabla]{babel}
\usepackage{csquotes}
\usepackage{geometry}
\geometry{a4paper,margin=2.5cm}

% --- Matemáticas y símbolos -------------------------------------------------
\usepackage{amsmath}
\usepackage{amssymb}
\usepackage{siunitx}
\sisetup{output-decimal-marker={,}, per-mode=symbol}

% --- Figuras, tablas, código -------------------------------------------------
\usepackage{graphicx}
\usepackage{booktabs}
\usepackage{caption}
\usepackage{subcaption}
\usepackage{listings}
\usepackage{xcolor}
\usepackage{enumitem}
\usepackage{float}

\lstdefinestyle{python}{
    language=Python,
    basicstyle=\ttfamily\footnotesize,
    keywordstyle=\color{blue!70!black}\bfseries,
    commentstyle=\color{gray!70!black}\itshape,
    stringstyle=\color{orange!70!black},
    numbers=left,
    numberstyle=\tiny\color{gray},
    breaklines=true,
    showstringspaces=false,
    frame=single,
    captionpos=b,
    tabsize=4,
}
\lstset{style=python}

% --- Bibliografía (biblatex, estilo IEEE) -----------------------------------
\usepackage[
    backend=biber,
    style=ieee,
    sorting=none,
    maxbibnames=99,
]{biblatex}
\addbibresource{references_U4.bib}

% --- Enlaces -----------------------------------------------------------------
\usepackage[colorlinks=true,linkcolor=blue!50!black,citecolor=blue!50!black,urlcolor=blue!50!black]{hyperref}
\usepackage[capitalise]{cleveref}

% --- Encabezados ---------------------------------------------------------
\usepackage{fancyhdr}
\pagestyle{fancy}
\setlength{\headheight}{14pt}
\fancyhf{}
\fancyhead[L]{GA-SUM-05 / PE-U4 --- Ley de Amdahl con Apache Spark}
\fancyhead[R]{Equipo ACC --- Aplicaciones Distribuidas (ISR-701)}
\fancyfoot[C]{\thepage}

\title{}
\author{}
\date{}

\begin{document}

% =============================================================================
    \begin{titlepage}
        \centering
        \vspace*{1cm}
        {\large\textsc{Universidad Técnica Estatal de Quevedo}}\\[0.2cm]
        {\normalsize Facultad de Ciencias de la Computación}\\
        {\normalsize Carrera de Ingeniería de Software}\\[1.5cm]

        {\small INSTRUMENTO DE EVALUACIÓN --- INFORME TÉCNICO}\\[0.3cm]
        \rule{\linewidth}{0.5pt}\\[0.4cm]
        {\Large\bfseries GA-SUM-05 / PE-U4 --- Práctica Experimental Unidad 4}\\[0.4cm]
        {\large Procesamiento Distribuido con Apache Spark: comprobación experimental de la Ley de
        Amdahl aplicada al Proyecto Fin de Curso}\\[0.4cm]
        \rule{\linewidth}{0.5pt}\\[1.2cm]

        \begin{tabular}{ll}
            \textbf{Asignatura:}          & Aplicaciones Distribuidas (ISR-701) \\
            \textbf{Unidad:}               & Unidad 4 --- Cómputo Paralelo y Distribuido \\
            \textbf{Período académico:}    & 2026--2027 PPA \\
            \textbf{Docente responsable:}  & Gleiston C. Guerrero-Ulloa, M.Sc. \\
        \end{tabular}

        \vspace{1.2cm}
        \textbf{Equipo:} D --- ACC - Soporte Técnico ISP\\[0.6cm]

        \begin{tabular}{ll}
            Alvarez Parraga Jeremy Alexis &
            \href{mailto:jalvarezp3@uteq.edu.ec}{\texttt{jalvarezp3@uteq.edu.ec}} \\

            Aucatoma Celorio Jhinson Stalyn &
            \href{mailto:jaucatomac@uteq.edu.ec}{\texttt{jaucatomac@uteq.edu.ec}} \\

            Carpio Mendoza Carlos Jose &
            \href{mailto:carlospatroner@gmail.com}{\texttt{carlospatroner@gmail.com}} \\
        \end{tabular}

        \vfill
        {\small Repositorio: \url{https://github.com/carlospatroner-boop/pe-u4-spark-Soporte-Tecnico-ISP}}\\[0.3cm]
    \end{titlepage}
\subsection{Contexto: el equipo y el dominio de referencia}
\label{sec:contexto-pfc}

El equipo ACC trabaja sobre el dominio \textbf{Soporte Técnico ISP}: la gestión de tickets de
soporte técnico de un proveedor de servicios de internet. Un ISP recibe, en operación normal, un
volumen de solicitudes de soporte varios órdenes de magnitud mayor al que un análisis
secuencial en un único hilo puede procesar con latencia razonable: clasificación de incidencias
por tipo, agregación por zona geográfica o estado, cruce con catálogos de referencia (clientes,
regiones, técnicos disponibles) y priorización automática de la cola de atención son operaciones
que, sobre históricos de cientos de miles o millones de tickets, se benefician directamente de un
motor de procesamiento distribuido. La decisión de negocio concreta que este trabajo busca
sustentar es la \textbf{priorización automática de la cola de atención técnica}: a partir de una
columna derivada de prioridad y de la agregación de tickets por estado/zona, un ISP real podría
decidir en qué zonas reforzar personal técnico y qué incidencias escalar primero, en lugar de
atenderlas estrictamente por orden de llegada.

No existe un histórico público de tickets internos de un ISP ecuatoriano al que el equipo tenga
acceso legítimo, por lo que se utilizó como sustituto un conjunto de datos real y público con
una estructura equivalente: \emph{CGB --- Consumer Complaints Data} de la Federal Communications
Commission de Estados Unidos, en el que cada fila es, literalmente, un \emph{ticket}
(\texttt{Ticket ID}) de queja de un consumidor contra un proveedor de telecomunicaciones,
clasificado por tipo de servicio (\emph{Phone}, \emph{Internet}, \emph{TV}), tipo de problema
específico, método de acceso y ubicación geográfica --la misma estructura de tipo de incidencia,
canal y zona que tendría el sistema de tickets interno de un ISP--. El dataset, su procedencia,
licencia y volumen se documentan en detalle en la \Cref{sec:procedimiento}.
\end{document}
